\documentclass[tikz,border=8pt]{standalone}
\usepackage{tikz}
\usepackage{amsmath}
\usepackage{amssymb}
\usepackage{pgfplots}
\pgfplotsset{compat=1.18}
\usepackage{ctex}
\usetikzlibrary{calc,positioning,arrows.meta,matrix}

% LC 1 Two Sum：哈希查找演化（4 步快照）
% 数组 [2,7,11,15], target=9

\begin{document}
\begin{tikzpicture}[
  every node/.style={font=\small},
  cell/.style={draw=gray!70, rectangle,
    minimum size=0.8cm, inner sep=0pt, thick, fill=white},
  curcell/.style={draw=orange!70, rectangle,
    minimum size=0.8cm, inner sep=0pt, very thick, fill=orange!10},
  foundcell/.style={draw=green!60!black, rectangle,
    minimum size=0.8cm, inner sep=0pt, very thick, fill=green!10},
  hashcell/.style={draw=blue!50, rectangle, rounded corners=2pt,
    minimum width=2.6cm, minimum height=0.65cm, thick, fill=blue!8},
  resultbox/.style={draw=green!60!black, fill=green!8, rounded corners=3pt,
    font=\small\bfseries, inner sep=5pt},
]

%% ── 标题 ──────────────────────────────────────────────────
\node[font=\large\bfseries] at (8.0, 1.2)
  {哈希查找（LC 1 Two Sum）};
\node[font=\small,gray] at (8.0, 0.55)
  {数组 \texttt{[2, 7, 11, 15]}，target = 9：遍历时用哈希表记录补数};

%% ── 顶部：原始数组 ──────────────────────────────────────
\node[font=\scriptsize,gray] at (2.0, -0.1) {数组：};
\foreach \i/\v in {0/2, 1/7, 2/11, 3/15} {
  \node[cell,font=\ttfamily] (arr\i) at (3.2+\i*1.0, -0.1) {\v};
  \node[font=\tiny,gray] at (3.2+\i*1.0, -0.65) {\i};
}

%% ══════════════════════════════════════════════════════════
%% 快照列 x 位置
\def\sA{1.6}
\def\sB{5.0}
\def\sC{8.8}
\def\sD{13.0}
\def\snapY{-1.8}

%% ── 快照 1：处理 nums[0]=2，need=9-2=7，dict 空，存入 {2:0} ──
\node[font=\scriptsize\bfseries,blue!70] at (\sA, \snapY) {步骤 1：idx=0，值=2};
\node[font=\tiny,gray] at (\sA, \snapY-0.4) {need=9-2=7，dict中无7};

% 当前元素高亮
\node[curcell,font=\ttfamily] at (3.2, -0.1) {2};
\node[font=\tiny,orange!70] at (3.2, -1.1) {$\uparrow$ 当前};

% dict 状态
\node[font=\scriptsize,anchor=east] at (\sA-0.3, \snapY-1.0) {dict:};
\node[hashcell,font=\ttfamily] at (\sA+0.9, \snapY-1.0) {2 $\to$ 0};
\node[font=\tiny,gray] at (\sA, \snapY-1.65) {存入 \{2: idx 0\}};

%% ── 快照 2：处理 nums[1]=7，need=9-7=2，dict中有 2 → 返回 [0,1] ──
\node[font=\scriptsize\bfseries,blue!70] at (\sB, \snapY) {步骤 2：idx=1，值=7};
\node[font=\tiny,gray] at (\sB, \snapY-0.4) {need=9-7=2，dict 中有 2!};

% 当前元素高亮
\node[foundcell,font=\ttfamily] at (4.2, -0.1) {7};
\node[font=\tiny,green!60!black] at (4.2, -1.1) {$\uparrow$ 命中};

% dict 状态（包含 2:0，加上 7:1）
\node[font=\scriptsize,anchor=east] at (\sB-0.5, \snapY-0.85) {dict:};
\node[hashcell,font=\ttfamily] at (\sB+0.7, \snapY-0.85) {2 $\to$ 0};
\node[hashcell,draw=green!60!black,fill=green!8,font=\ttfamily]
  at (\sB+0.7, \snapY-1.55) {7 $\to$ 1};
\node[font=\tiny,green!60!black] at (\sB+0.7, \snapY-2.1) {查到 2 在 idx=0};

%% ── 快照 3：步骤 3（假设继续）nums[2]=11，need=-2，无 ──
\node[font=\scriptsize\bfseries,blue!70] at (\sC, \snapY) {步骤 3（若继续）};
\node[font=\tiny,gray] at (\sC, \snapY-0.4) {idx=2，值=11，need=-2，无};

\node[curcell,font=\ttfamily] at (5.2, -0.1) {11};
\node[font=\tiny,orange!70] at (5.2, -1.1) {$\uparrow$ 当前};

\node[font=\scriptsize,anchor=east] at (\sC-0.5, \snapY-0.85) {dict:};
\node[hashcell,font=\ttfamily] at (\sC+0.7, \snapY-0.85) {2 $\to$ 0};
\node[hashcell,font=\ttfamily] at (\sC+0.7, \snapY-1.55) {7 $\to$ 1};
\node[hashcell,font=\ttfamily] at (\sC+0.7, \snapY-2.25) {11 $\to$ 2};
\node[font=\tiny,gray] at (\sC+0.7, \snapY-2.8) {存入 \{11: idx 2\}};

%% ── 快照 4：最终结果 ──
\node[font=\scriptsize\bfseries,blue!70] at (\sD, \snapY) {最终结果};
\node[font=\tiny,gray] at (\sD, \snapY-0.4) {步骤 2 即命中，提前返回};

% 高亮 idx 0 和 idx 1
\node[foundcell,font=\ttfamily] at (3.2, -0.1) {2};
\node[foundcell,font=\ttfamily] at (4.2, -0.1) {7};

\node[resultbox] at (\sD, \snapY-1.4)
  {\textcolor{green!50!black}{返回 [0, 1]}};

\node[font=\scriptsize,draw=gray!40,fill=white,thin,rounded corners=2pt,
      inner sep=4pt,align=center]
  at (\sD, \snapY-2.5)
  {nums[0]+nums[1]\\=2+7=9=target};

%% ── 分隔线 ──────────────────────────────────────────────
\draw[gray!25,dashed] (0.0, -5.2) -- (16.0, -5.2);

%% ── 算法说明 ─────────────────────────────────────────────
\node[draw=gray!50,fill=white,rounded corners=3pt,font=\small,
      inner sep=8pt,align=left]
  at (8.0, -6.2)
  {\textbf{Two Sum 哈希算法（O(n) 时间，O(n) 空间）：}\\[4pt]
   遍历数组，对每个元素 \texttt{x}：\\[2pt]
   \quad 1. 计算补数 \texttt{need = target - x}\\[2pt]
   \quad 2. 若 \texttt{need} 在 dict 中 $\to$ 找到答案，返回 \texttt{[dict[need], i]}\\[2pt]
   \quad 3. 否则将 \texttt{x} $\to$ \texttt{i} 存入 dict，继续下一元素};

\end{tikzpicture}
\end{document}
